\section{本讲小结}

网络简化应先保持连接关系，再利用当前端口的边界条件判断等势。保持两端分别不动的对称变换给出对应点等势；交换两端的变换通常给出电势互补。删去无流边或合并等势点，都要说明其条件，更换端口后应重新判断。

立方体与高维超立方体可用等势层降低未知量，跨层边数决定并联电阻。完全图则利用内部节点置换得到相同中间电势。不能完全简化的部分继续列节点方程，保留物理约束比追求图形上的对称更重要。

星三角变换保留三个端口。使用时应标清对应端点，优先选择能立即形成串并联的局部结构，求内部电流时回到真实网络。最后用恢复原阻值、短路、断路、增删边单调性及节点守恒检查所得电阻。

